\documentclass[tikz,border=0pt]{standalone}
\usepackage{fontspec}
\usepackage{xeCJK}
\usepackage{xcolor}
\usepackage{tikz}
\usetikzlibrary{positioning,calc,shadows.blur,arrows.meta,fit,backgrounds}
\usepackage{microtype}
\usepackage[strings]{underscore}
\setlength{\parindent}{0pt}

% ============================================================
% Fonts
% ============================================================
\IfFontExistsTF{Microsoft YaHei}{
  \setCJKmainfont{Microsoft YaHei}
  \setCJKsansfont{Microsoft YaHei}
  \setCJKmonofont{Microsoft YaHei}
}{
  \setCJKmainfont{Noto Sans CJK SC}
  \setCJKsansfont{Noto Sans CJK SC}
  \setCJKmonofont{Noto Sans Mono CJK SC}
}
\setmonofont{DejaVu Sans Mono}
\setsansfont{DejaVu Sans}
\renewcommand\familydefault{\sfdefault}

% ============================================================
% Colors
% ============================================================
\definecolor{Navy}{HTML}{173B7A}
\definecolor{Blue}{HTML}{2878D4}
\definecolor{BlueSoft}{HTML}{EEF6FF}
\definecolor{Teal}{HTML}{0F9D87}
\definecolor{TealSoft}{HTML}{EEFDF9}
\definecolor{Orange}{HTML}{F28C18}
\definecolor{OrangeSoft}{HTML}{FFF7EC}
\definecolor{Rose}{HTML}{E84A64}
\definecolor{RoseSoft}{HTML}{FFF0F3}
\definecolor{Purple}{HTML}{6C55C9}
\definecolor{PurpleSoft}{HTML}{F4F1FF}
\definecolor{Green}{HTML}{28A65A}
\definecolor{GreenSoft}{HTML}{EEF9F1}
\definecolor{Text}{HTML}{24364B}
\definecolor{Muted}{HTML}{6E7B8E}
\definecolor{CodeBg}{HTML}{F5F7FA}
\definecolor{CodeBorder}{HTML}{D8E0EA}
\definecolor{Line}{HTML}{D8E2ED}
\definecolor{Gold}{HTML}{F3B63A}

% ============================================================
% Layout parameters -- maintain the geometry here.
% Almost every placement below is derived from these values.
% ============================================================
\newlength{\PageW}       \setlength{\PageW}{21cm}
\newlength{\PageH}       \setlength{\PageH}{29.7cm}
\newlength{\MainW}       \setlength{\MainW}{18.6cm}
\newlength{\HalfW}       \setlength{\HalfW}{9.15cm}
\newlength{\PanelGap}    \setlength{\PanelGap}{3mm}
\newlength{\RowGap}      \setlength{\RowGap}{3mm}
\newlength{\PanelH}      \setlength{\PanelH}{6.22cm}
\newlength{\WidePanelH}  \setlength{\WidePanelH}{4.06cm}
\newlength{\BottomH}     \setlength{\BottomH}{5.10cm}
\newlength{\LadderW}     \setlength{\LadderW}{10.1cm}
\newlength{\TopTenW}     \setlength{\TopTenW}{8.2cm}

\newlength{\PanelPad}    \setlength{\PanelPad}{1.8mm}
\newlength{\CardsTop}    \setlength{\CardsTop}{8.3mm}
\newlength{\CardW}       \setlength{\CardW}{4.34cm}
\newlength{\CardH}       \setlength{\CardH}{2.52cm}
\newlength{\CardGapX}    \setlength{\CardGapX}{1.1mm}
\newlength{\CardGapY}    \setlength{\CardGapY}{0.8mm}
\newlength{\WideCardW}   \setlength{\WideCardW}{4.42cm}
\newlength{\WideCardH}   \setlength{\WideCardH}{2.90cm}
\newlength{\WideCardGap} \setlength{\WideCardGap}{1.8mm}

% ============================================================
% Global TikZ styles
% ============================================================
\tikzset{
  panel/.style={
    rounded corners=3.2mm,
    line width=0.7pt,
    fill=white,
    blur shadow={shadow blur steps=5,shadow xshift=0.6mm,shadow yshift=-0.6mm,shadow opacity=10}
  },
  header/.style={
    rounded corners=2.6mm,
    minimum height=9mm,
    text=white,
    font=\bfseries\large,
    inner xsep=3.2mm,
    align=left
  },
  itemcard/.style={
    rounded corners=2.2mm,
    draw=Line,
    fill=white,
    line width=0.45pt,
    inner sep=2.1mm,
    align=left
  },
  num/.style={
    circle,
    minimum size=5.8mm,
    inner sep=0pt,
    text=white,
    font=\bfseries\scriptsize
  },
  code/.style={
    rounded corners=1.5mm,
    draw=CodeBorder,
    fill=CodeBg,
    line width=0.35pt,
    inner sep=1.5mm,
    align=left,
    font=\ttfamily\fontsize{6.0}{7.0}\selectfont,
    text=Text
  },
  arrow/.style={-{Latex[length=2.1mm,width=1.7mm]},line width=0.8pt}
}

% ============================================================
% Reusable components
% ============================================================
% Panel: name, placement, width, height, color, title, tagline
\newcommand{\MakePanel}[7]{%
  \node[panel,draw=#5!85,minimum width=#3,minimum height=#4,anchor=north west] (#1) at #2 {};
  % Draw the header as an empty bar, then place title/tagline independently.
  % This prevents a long tagline from colliding with the title.
  \node[header,fill=#5,minimum width=\dimexpr#3-1.5mm\relax,anchor=north west]
    (#1-head) at ([xshift=.75mm,yshift=-.75mm]#1.north west) {};
  \node[anchor=west,text=white,font=\bfseries\large,
        text width=\dimexpr#3-36mm\relax,align=left]
    at ([xshift=3.2mm]#1-head.west) {#6};
  \node[anchor=east,text=white!78,font=\fontsize{6.0}{6.8}\selectfont,
        text width=30mm,align=right]
    at ([xshift=-3.2mm]#1-head.east) {#7};
}

% Standard card: name, placement, number, color, title, body, code, title-font-size
\newcommand{\RefCard}[8]{%
  \node[itemcard,minimum width=\CardW,minimum height=\CardH,anchor=north west] (#1) at #2 {};
  \node[num,fill=#4,anchor=north west] at ([xshift=1.7mm,yshift=-1.7mm]#1.north west) {#3};
  \node[anchor=north west,text=#4!82!black,font=\bfseries\fontsize{#8}{8.7}\selectfont,
        text width=\dimexpr\CardW-12mm\relax]
    at ([xshift=8.7mm,yshift=-1.8mm]#1.north west) {#5};
  \node[anchor=north west,text=Text,font=\fontsize{6.75}{7.9}\selectfont,
        text width=\dimexpr\CardW-4.4mm\relax]
    at ([xshift=2mm,yshift=-8mm]#1.north west) {#6};
  \node[code,anchor=north west,text width=\dimexpr\CardW-4.7mm\relax]
    at ([xshift=2mm,yshift=-15.3mm]#1.north west) {#7};
}

% Wide-row card used in section E.
\newcommand{\WideRefCard}[8]{%
  \node[itemcard,minimum width=\WideCardW,minimum height=\WideCardH,anchor=north west] (#1) at #2 {};
  \node[num,fill=#4,anchor=north west] at ([xshift=1.7mm,yshift=-1.7mm]#1.north west) {#3};
  \node[anchor=north west,text=#4!80!black,font=\bfseries\fontsize{#8}{7.6}\selectfont,
        text width=\dimexpr\WideCardW-10mm\relax]
    at ([xshift=8.7mm,yshift=-1.8mm]#1.north west) {#5};
  \node[anchor=north west,text=Text,font=\fontsize{6.55}{7.7}\selectfont,
        text width=\dimexpr\WideCardW-4.4mm\relax]
    at ([xshift=2mm,yshift=-9mm]#1.north west) {#6};
  \node[code,anchor=north west,text width=\dimexpr\WideCardW-4.7mm\relax]
    at ([xshift=2mm,yshift=-18.3mm]#1.north west) {#7};
}

% Relative placement helpers.
\newcommand{\FirstCardPos}[1]{([xshift=\PanelPad,yshift=-\CardsTop]#1.north west)}
\newcommand{\RightCardPos}[1]{([xshift=\CardGapX]#1.north east)}
\newcommand{\BelowCardPos}[1]{([yshift=-\CardGapY]#1.south west)}
\newcommand{\WideFirstCardPos}[1]{([xshift=\PanelPad,yshift=-\CardsTop]#1.north west)}
\newcommand{\WideRightCardPos}[1]{([xshift=\WideCardGap]#1.north east)}

\begin{document}
\begin{tikzpicture}[x=1cm,y=1cm]
  % A4 page. This is the only intentionally absolute geometry.
  \path[use as bounding box] (0,0) rectangle (21,-29.7);

  % Master anchors. Change MainW to resize/recenter the entire card.
  \coordinate (PageTop) at (10.5,-0.7);
  \coordinate (MainNW) at ($(PageTop)+(-.5\MainW,0)$);

  % ============================================================
  % Header -- chained vertically, no absolute Y coordinates.
  % ============================================================
  \node[anchor=north,text=Navy,font=\bfseries\fontsize{20}{22}\selectfont,
        align=center,text width=\MainW] (title)
    at (PageTop) {C 语言：把“靠人不出错”变成“靠编译器不出错”};

  \node[below=1.3mm of title.south,anchor=north,text=Muted,
        font=\bfseries\fontsize{10.8}{12}\selectfont] (subtitle)
    {20 条可落地实践 · Reference Card};

  \node[below=3.6mm of subtitle.south,anchor=north,rounded corners=2.2mm,
        draw=Gold!85!black,fill=Gold!8,line width=0.55pt,
        minimum width=17.1cm,minimum height=0.95cm] (principle) {};
  \node[anchor=west,text=Gold!65!black,font=\bfseries\large]
    at ([xshift=2.8mm]principle.west) {\Large\textbf{!}};
  \node[anchor=west,text=Navy,font=\bfseries\fontsize{7.2}{8.6}\selectfont,text width=15.6cm]
    at ([xshift=8mm]principle.west)
    {原则：能让类型系统检查的，不靠代码评审；能让编译器报错的，不靠运行测试；能在编译期失败的，不留到运行期。};

  % ============================================================
  % Row 1: A + B
  % B follows A; moving/resizing A moves B automatically.
  % ============================================================
  \MakePanel{A}{([xshift=-.5\MainW,yshift=-3mm]principle.south)}{\HalfW}{\PanelH}{Blue}
    {A. 编译期断言与 API 契约}{让约束前置到编译期}
  \MakePanel{B}{([xshift=\PanelGap]A.north east)}{\HalfW}{\PanelH}{Teal}
    {B. 类型系统约束}{用类型表达意图，让错误无处遁形}

  \RefCard{A1}{\FirstCardPos{A}}{1}{Blue}{_Static_assert()}
    {把结构大小、偏移、位图容量等假设写成编译期断言。}
    {_Static_assert(sizeof(struct\\packet)==8, "layout changed");}{8.0}
  \RefCard{A2}{\RightCardPos{A1}}{2}{Blue}{warning / error}
    {旧接口可 warning；绝对禁止的接口直接 error。}
    {__attribute__((warning("...")))\\__attribute__((error("...")))}{7.3}
  \RefCard{A3}{\BelowCardPos{A1}}{3}{Blue}{deprecated}
    {渐进式重构：先 deprecated，再 error，最后删除。}
    {__attribute__((deprecated("use new_api")))}{8.0}
  \RefCard{A4}{\RightCardPos{A3}}{4}{Blue}{warn_unused_result / [[nodiscard]]}
    {错误码不允许被静默忽略。}
    {int flash_program(...)\\__attribute__((warn_unused_result));}{7.1}

  \RefCard{B5}{\FirstCardPos{B}}{5}{Teal}{format(printf, ...)}
    {让自定义日志函数获得 printf 格式检查。}
    {void log_msg(const char *fmt, ...)\\__attribute__((format(printf,1,2)));}{7.7}
  \RefCard{B6}{\RightCardPos{B5}}{6}{Teal}{nonnull}
    {把“参数不可为 NULL”写成编译器可理解的契约。}
    {int pt_execute(...)\\__attribute__((nonnull(1,2)));}{8.0}
  \RefCard{B7}{\BelowCardPos{B5}}{7}{Teal}{const}
    {不是风格问题，而是“此函数没有资格修改它”。}
    {int pt_execute(const struct\\pt_command *cmd, struct pt_context *ctx);}{8.0}
  \RefCard{B8}{\RightCardPos{B7}}{8}{Teal}{不同类型表达不同概念}
    {不要所有东西都用 int；必要时用 enum / typedef struct 封装强类型。}
    {typedef struct { unsigned value; }\\pt_port_t;}{7.5}

  % ============================================================
  % Row 2: C + D. C follows A vertically; D follows C horizontally.
  % ============================================================
  \MakePanel{C}{([yshift=-\RowGap]A.south west)}{\HalfW}{\PanelH}{Orange}
    {C. 封装与可见性}{隐藏实现，暴露契约}
  \MakePanel{D}{([xshift=\PanelGap]C.north east)}{\HalfW}{\PanelH}{Rose}
    {D. 泛型、表驱动与初始化}{让数据描述变化，让编译器检查一致性}

  \RefCard{C9}{\FirstCardPos{C}}{9}{Orange}{opaque type}
    {公共头文件只暴露 struct 声明，不暴露内部字段。}
    {struct pt_device;}{8.0}
  \RefCard{C10}{\RightCardPos{C9}}{10}{Orange}{static 限定作用域}
    {不需要导出的函数，一律 static，缩小可访问面。}
    {static int helper_do_work(void);\\/* only this translation unit */}{7.6}
  \RefCard{C11}{\BelowCardPos{C9}}{11}{Orange}{switch(enum) + -Wswitch-enum}
    {新增状态后，遗漏分支处理时让编译器报警/报错。}
    {switch (state) {\\case PT_IDLE: ...; break;\\case PT_BUSY: ...; break;\\}}{7.0}
  \RefCard{C12}{\RightCardPos{C11}}{12}{Orange}{函数指针类型约束回调}
    {用 typedef 的函数签名约束表驱动回调接口。}
    {typedef int (*pt_handler_t)(\\const struct pt_command *,\\struct pt_context *);}{7.3}

  \RefCard{D13}{\FirstCardPos{D}}{13}{Rose}{_Generic}
    {做有限的编译期类型分派，限制可接受的参数类型。}
    {\#define pt_abs(x) _Generic((x),\\int: pt_abs_int, long: pt_abs_long)(x)}{8.0}
  \RefCard{D14}{\RightCardPos{D13}}{14}{Rose}{ARRAY_SIZE + _Static_assert}
    {检查 handler table / opcode table 的完整性。}
    {_Static_assert(ARRAY_SIZE(handlers)\\== PT_OPCODE_COUNT, "table mismatch");}{7.0}
  \RefCard{D15}{\BelowCardPos{D13}}{15}{Rose}{designated initializer}
    {减少位置填错；结构体增删字段时更稳。}
    {struct ops x = {\\.open = f_open,\\.read = f_read,\\};}{7.2}
  \RefCard{D16}{\RightCardPos{D15}}{16}{Rose}{表驱动设计}
    {把变化项放进表，把流程代码变成受约束的数据。}
    {static const struct pt_op ops[] = {\\{}[OP_OPEN] = { .fn = f_open },\\{}[OP_READ] = { .fn = f_read },\\};}{8.0}

  % ============================================================
  % Row 3: E follows C/D. Four cards are chained left-to-right.
  % ============================================================
  \MakePanel{E}{([yshift=-\RowGap]C.south west)}{\MainW}{\WidePanelH}{Purple}
    {E. 编译器开关与工程纪律}{打开编译器的“全量体检模式”}

  \WideRefCard{E17}{\WideFirstCardPos{E}}{17}{Purple}{-Werror}
    {warning 只有变成 error，才真正形成约束。}
    {-Wall -Wextra -Werror}{6.9}
  \WideRefCard{E18}{\WideRightCardPos{E17}}{18}{Purple}{-Wmissing-prototypes / -Wimplicit-fallthrough}
    {强制 API 走头文件；防止 switch 漏 break。}
    {-Wmissing-prototypes\\-Wimplicit-fallthrough}{6.9}
  \WideRefCard{E19}{\WideRightCardPos{E18}}{19}{Purple}{不要滥用 cast}
    {cast 往往是在“关闭编译器”，应视为需要解释的危险动作。}
    {/* bad */ int x = (int)ptr;\\/* better: fix the type design */}{6.9}
  \WideRefCard{E20}{\WideRightCardPos{E19}}{20}{Purple}{__must_check / BUILD_BUG_ON()}
    {把常用约束包装成统一宏，形成工程规范。}
    {\#define BUILD_BUG_ON(c)\\_Static_assert(!(c), "BUILD_BUG_ON")}{6.9}

  % ============================================================
  % Bottom row: ladder + Top 10, both relative to E.
  % ============================================================
  \MakePanel{L}{([yshift=-\RowGap]E.south west)}{\LadderW}{\BottomH}{Blue}
    {约束等级（从弱到强）}{}
  \MakePanel{T}{([xshift=\PanelGap]L.north east)}{\TopTenW}{\BottomH}{Green}
    {底层 C 项目优先推荐}{Top 10}

  % Ladder starts from L's header; every later box is positioned from its predecessor.
  \node[rounded corners=1.5mm,draw=Line,fill=white,minimum width=1.18cm,minimum height=.78cm,
        font=\bfseries\fontsize{5.8}{6.3}\selectfont,anchor=north west]
        (l1) at ([xshift=4.2mm,yshift=-11.4mm]L.north west) {注释约束};
  \node[right=1.5mm of l1,rounded corners=1.5mm,draw=Line,fill=BlueSoft,minimum width=1.15cm,minimum height=.78cm,
        font=\bfseries\fontsize{5.8}{6.3}\selectfont] (l2) {warning};
  \node[right=1.5mm of l2,rounded corners=1.5mm,draw=Line,fill=GreenSoft,minimum width=1.25cm,minimum height=.78cm,
        font=\bfseries\fontsize{5.6}{6.1}\selectfont,align=center] (l3) {warning\\$\to$error};
  \node[right=1.5mm of l3,rounded corners=1.5mm,draw=Line,fill=OrangeSoft,minimum width=1.32cm,minimum height=.78cm,
        font=\bfseries\fontsize{5.5}{6.1}\selectfont,align=center] (l4) {编译期\\断言};
  \node[right=1.5mm of l4,rounded corners=1.5mm,draw=Line,fill=RoseSoft,minimum width=1.22cm,minimum height=.78cm,
        font=\bfseries\fontsize{5.6}{6.1}\selectfont] (l5) {类型系统};
  \node[right=1.5mm of l5,rounded corners=1.5mm,draw=Line,fill=PurpleSoft,minimum width=1.22cm,minimum height=.78cm,
        font=\bfseries\fontsize{5.6}{6.1}\selectfont,align=center] (l6) {API\\不可见};
  \node[right=1.5mm of l6,rounded corners=1.5mm,draw=Line,fill=TealSoft,minimum width=1.44cm,minimum height=.78cm,
        font=\bfseries\fontsize{5.3}{5.9}\selectfont,align=center] (l7) {错误操作\\无法表达};

  \foreach \a/\b in {l1/l2,l2/l3,l3/l4,l4/l5,l5/l6,l6/l7}
    \draw[arrow,draw=Muted] (\a.east) -- (\b.west);

  \node[below=1.3mm of l1.south,text=Muted,font=\fontsize{5.5}{6.5}\selectfont,align=center,text width=1.35cm] {仅靠人工\\容易遗漏};
  \node[below=1.3mm of l2.south,text=Muted,font=\fontsize{5.5}{6.5}\selectfont,align=center,text width=1.35cm] {编译器提示\\但不阻止};
  \node[below=1.3mm of l3.south,text=Muted,font=\fontsize{5.5}{6.5}\selectfont,align=center,text width=1.35cm] {将警告变为\\强制约束};
  \node[below=1.3mm of l4.south,text=Muted,font=\fontsize{5.5}{6.5}\selectfont,align=center,text width=1.35cm] {在编译期\\发现问题};
  \node[below=1.3mm of l5.south,text=Muted,font=\fontsize{5.5}{6.5}\selectfont,align=center,text width=1.35cm] {用类型表达\\约束};
  \node[below=1.3mm of l6.south,text=Muted,font=\fontsize{5.5}{6.5}\selectfont,align=center,text width=1.35cm] {从接口层面\\防止误用};
  \node[below=1.3mm of l7.south,text=Muted,font=\fontsize{5.5}{6.5}\selectfont,align=center,text width=1.45cm] {从根本上\\不允许出错};

  \draw[-{Latex[length=2.5mm,width=2mm]},line width=1.6pt,draw=Blue!80]
    ([xshift=4.2mm,yshift=7.2mm]L.south west) -- ([xshift=-4.2mm,yshift=7.2mm]L.south east);
  \node[anchor=north,text=Navy,font=\bfseries\fontsize{7.0}{8.0}\selectfont]
    at ([yshift=5.8mm]L.south) {更强的约束  $\longrightarrow$  更少的 Bug  $\longrightarrow$  更可靠的系统};

  % Top-10 list: two relative columns, each item follows the one above it.
  \node[num,fill=Green,minimum size=4.4mm,font=\bfseries\fontsize{5.8}{6.1}\selectfont,anchor=north west]
    (t1n) at ([xshift=4.4mm,yshift=-10.6mm]T.north west) {1};
  \node[anchor=west,text=Text,font=\fontsize{6.6}{7.4}\selectfont] (t1) at ([xshift=3.4mm]t1n.east) {const};
  \foreach \prev/\num/\txt in {t1n/2/_Static_assert,t2n/3/warn_unused_result,t3n/4/{format(printf,...)},t4n/5/{deprecated / error}}{
    \pgfmathtruncatemacro{\thisnum}{\num}
    \node[below=1.3mm of \prev,num,fill=Green,minimum size=4.4mm,font=\bfseries\fontsize{5.8}{6.1}\selectfont,anchor=north] (t\thisnum n) {\thisnum};
    \node[anchor=west,text=Text,font=\fontsize{6.6}{7.4}\selectfont] at ([xshift=3.4mm]t\thisnum n.east) {\txt};
  }

  \node[num,fill=Green,minimum size=4.4mm,font=\bfseries\fontsize{5.8}{6.1}\selectfont,anchor=north west]
    (t6n) at ([xshift=2.0mm,yshift=-10.6mm]T.north) {6};
  \node[anchor=west,text=Text,font=\fontsize{6.35}{7.2}\selectfont] at ([xshift=3.4mm]t6n.east) {-Wswitch-enum};
  \foreach \prev/\num/\txt in {t6n/7/-Wmissing-prototypes,t7n/8/{static + opaque struct},t8n/9/{强类型 typedef struct},t9n/10/{designated initializer + 表驱动}}{
    \pgfmathtruncatemacro{\thisnum}{\num}
    \node[below=1.3mm of \prev,num,fill=Green,minimum size=4.4mm,font=\bfseries\fontsize{5.8}{6.1}\selectfont,anchor=north] (t\thisnum n) {\thisnum};
    \node[anchor=west,text=Text,font=\fontsize{6.35}{7.2}\selectfont] at ([xshift=3.4mm]t\thisnum n.east) {\txt};
  }

  \node[rounded corners=1.7mm,draw=Green!25,fill=GreenSoft,
        minimum width=\dimexpr\TopTenW-6mm\relax,minimum height=.72cm,anchor=south]
    (tip) at ([yshift=3.7mm]T.south) {};
  \node[anchor=west,text=Green!55!black,font=\bfseries\fontsize{6.5}{7.3}\selectfont]
    at ([xshift=2.8mm]tip.west) {小步坚持，长期受益。好的 C，源于更严谨的工程约束。};

  % Footer follows the bottom row instead of using a page coordinate.
  \node[below=4.0mm of T.south east,anchor=north east,text=Muted,font=\fontsize{5.5}{6.2}\selectfont]
    {XeLaTeX + TikZ · compiler-enforced C design};

\end{tikzpicture}
\end{document}
